\documentclass[11pt]{article}
\usepackage[T1]{fontenc}
\usepackage{lmodern,amsmath,amssymb,amsthm,mathtools,bm}
\usepackage[margin=1in]{geometry}
\usepackage{microtype,booktabs,longtable,array,enumitem}
\usepackage[hidelinks]{hyperref}
\usepackage{fancyvrb}
\setlength{\parindent}{0pt}
\setlength{\parskip}{5pt}
\setlist{nosep,leftmargin=*}
\newtheorem{theorem}{Theorem}[section]
\newtheorem{lemma}[theorem]{Lemma}
\newtheorem{proposition}[theorem]{Proposition}
\newtheorem{corollary}[theorem]{Corollary}
\theoremstyle{definition}\newtheorem{definition}[theorem]{Definition}
\theoremstyle{remark}\newtheorem{remark}[theorem]{Remark}
\newcommand{\C}{\mathbb C}
\newcommand{\Q}{\mathbb Q}
\newcommand{\PP}{\mathbb P}
\newcommand{\br}{\underline{\mathrm R}}
\newcommand{\crk}{\mathrm{cR}}
\newcommand{\Span}{\operatorname{span}}
\newcommand{\Seg}{\operatorname{Seg}}
\newcommand{\Sym}{\operatorname{Sym}}
\newcommand{\Gr}{\operatorname{Gr}}
\newcommand{\rank}{\operatorname{rank}}
\newcommand{\codim}{\operatorname{codim}}
\newcommand{\im}{\operatorname{im}}
\newcommand{\calS}{\mathcal S}
\newcommand{\calG}{\mathcal G}
\newcommand{\ann}{\operatorname{ann}}
\newcommand{\src}[2]{\footnote{#1. See source(s) [#2] in the Sources section.}}
\begin{document}
\begin{center}
{\LARGE\bfseries AC-03: Pair-space reduction and cubic obstructions}
\end{center}

\textbf{Status.} This report does not determine the exact complex border rank of
$3\times3$ matrix multiplication. Its numerical conclusion remains
\[
\boxed{17\leq\br(M_3)\leq20.}
\]
The two endpoints are published results, not new bounds. The retained prior
package verifies them and provides an exhaustive rank-17 reduction. This report
builds on that reduction with an equivalent three-factor target, exact scheme
sections, explicit grouped degenerations, and characteristic-zero obstructions
to two substantial classes of cubic-order constructions.

The reduced target is a $17\times9\times9$ tensor $F_E$ with 35 nonzero coordinates.
Its border rank is 17 if and only if $\br(M_3)=17$. Its cactus rank is at least 18,
but this is \emph{not} a border-rank exclusion. An explicit control tensor in the
same format has border rank 17 and cactus rank at least 18, illustrating the gap.

For the four-factor target inherited from the preceding report, all three
versions obtained by grouping two of its original matrix factors have exact
border rank 17. The displayed degenerations do not split into four-factor
rank-one terms. On the exclusion side, a 154\,350-column exterior-contraction map
rules out cubic-order constructions at all 24 normalized root-of-unity base
configurations. A separate integral-kernel and quadratic-ideal certificate rules
out the grid-plus-barycenter configuration. An independent 24-pattern calculation
excludes a restricted monomial Fourier ansatz for the pair plane.

Every restriction in those statements is essential. Arbitrary limiting
configurations, nonreduced limiting schemes, and higher-order degenerations have
not been excluded. No matching construction with 17, 18, or 19 terms is supplied.

\newpage
\tableofcontents
\newpage

\section{Problem, conventions, and the inherited reduction}

The repository asks for the exact border rank over $\C$ of
\begin{equation}\label{eq:M}
 M_3=\sum_{i,j,k=1}^{3}x_{ij}\otimes y_{jk}\otimes z_{ki}
 \in A\otimes B\otimes C,\qquad \dim A=\dim B=\dim C=9.
\end{equation}
The third factor uses the trace convention: contraction gives
$\operatorname{tr}(XYZ)$. Border rank permits limits of sums of decomposable
tensors; it is not the ordinary tensor rank asked for in AC-02. The public
AC-03 statement continues to list the exact value as unresolved.\src{OpenProblemsInNLA,
AC-03, statement retrieved 14 September 2026}{1}

The lower bound 17 is due to Conner--Harper--Landsberg. The upper certificate
retained here is Smirnov's length-20 degeneration.\src{Conner--Harper--Landsberg,
\emph{New lower bounds for matrix multiplication and the $3\times3$ determinant};
Smirnov, \emph{The bilinear complexity and practical algorithms for matrix
multiplication}}{2, 3}
The previous packages reproduce those bounds by exact arithmetic. The upper
certificate has constant tensor $M_3$, no negative Laurent residual, and error
at most $111|t|$ in Frobenius norm for $0<|t|\leq1$. The lower certificate
eliminates all candidate Borel-fixed pair data at rank 16.

For a projective variety $V\subset\PP W$, define
\begin{equation}\label{eq:grass}
\calG_r(V)=\overline{\{\Span(p_1,\ldots,p_r):p_i\in V,
                      \dim\Span(p_1,\ldots,p_r)=r\}}
             \subset\Gr(r,W).
\end{equation}
All dimensions of linear spaces in this report are vector-space dimensions.
Membership in \eqref{eq:grass} concerns a limit of \emph{entire spans}. It does
not require the limiting space to be spanned by a length-$r$ scheme contained
in that space.

Put
\begin{align}\label{eq:core}
 A'&=\Span(x_{21},x_{31}),&
 B'&=\Span(y_{21},y_{22},y_{31},y_{32}),&
 C'&=\Span(z_{31},z_{32}),\notag\\
 L&=A'\otimes B'\otimes C',&
 \calS&=\Span(M_3)\oplus L,&
 X&=\Seg(\PP A\times\PP B\times\PP C).
\end{align}
The 16 coordinate tensors of $L$ have disjoint support from $M_3$, so
$\dim\calS=17$. Let $D$ have basis $d_0$ and $d_{abc}$ indexed by those 16
coordinate tensors. Define
\begin{equation}\label{eq:Q}
 Q=d_0\otimes M_3+\sum_{(a,b,c)\in\mathcal B_L}
                         d_{abc}\otimes a\otimes b\otimes c.
\end{equation}
It has 43 nonzero unit coordinates and is concise in $D$.

\subsection{Audited input from the preceding package}

The preceding report proves the following computer-assisted statement. Its
complete verifier and unchanged archive are included under \texttt{prior/}.

\begin{theorem}[Inherited single-plane reduction]\label{thm:prior}
One has
\[
 \br(M_3)=17\quad\Longleftrightarrow\quad
 \calS\in\calG_{17}(X)\quad\Longleftrightarrow\quad
 \br_{D,A,B,C}(Q)=17.
\]
Moreover, any Borel-fixed simultaneous span data of a rank-17 degeneration,
after a cyclic rotation of $A,B,C$, have pair spaces
\begin{equation}\label{eq:canonicalpairs}
 E_{AB}=M_3(C^*)\oplus(A'\otimes B'),\qquad
 E_{BC}=M_3(A^*)\oplus(B'\otimes C').
\end{equation}
\end{theorem}

The classification starts from 343 Borel pair-space charts, including 29
one-parameter charts. The pair tests and square test leave 55 fixed spaces and
no parameter families. All $55^3=166\,375$ ordered triples are covered. Exactly
87 pass the three-factor test; they form 29 cyclic orbits, and each surviving
intersection is 17-dimensional. After cyclic rotation, each is exactly $\calS$.
The square threshold uses first jets at 17 general points; its modern framework
is discussed by Ma\'ndziuk.\src{Ma\'ndziuk, \emph{Border rank lower bounds beyond weak
border apolarity}, submitted 10 September 2026}{4}

The necessity in Theorem~\ref{thm:prior} is not merely an apolarity assertion.
The proof retains actual pair and triple spans in a projective incidence
closure, imposes $M_3$ in the triple span, and applies the Borel fixed-point
theorem. Thus the resulting triple space remains in $\calG_{17}(X)$. The
square and pair conditions are closed necessary conditions on those same data.
The complete enumeration forces the triple space to be $\calS$.

Conversely, a local basis of a family of 17-dimensional spans converging to
$\calS$ lifts $M_3$ to tensors of rank at most 17. The first-factor flattening of
$Q$ gives the last equivalence. Cyclic rotation is used to identify equivalent
ordered triples, not to require equal pair spaces.

\section{An equivalent 35-coordinate tensor}

Let
\begin{equation}\label{eq:E}
 E=\Span\left\{s_{ik}=\sum_{j=1}^3x_{ij}\otimes y_{jk}:1\leq i,k\leq3\right\}
       \oplus(A'\otimes B')\subset A\otimes B.
\end{equation}
Take a 17-dimensional space $D_E$ with basis $e_{ik}$ and $e_{ab}$ indexed by the
nine slices and the eight core coordinate tensors. Its inclusion tensor is
\begin{equation}\label{eq:FE}
 F_E=\sum_{i,j,k=1}^3e_{ik}\otimes x_{ij}\otimes y_{jk}
       +\sum_{a\in\mathcal B_{A'},\ b\in\mathcal B_{B'}}e_{ab}\otimes a\otimes b.
\end{equation}
The supports of the two sums are disjoint. Thus $F_E$ has 35 nonzero unit
coordinates and is concise in $D_E$.

\begin{theorem}[Single-pair equivalence]\label{thm:pair}
With $Z=\Seg(\PP A\times\PP B)$,
\[
\boxed{\br(M_3)=17\ \Longleftrightarrow\ E\in\calG_{17}(Z)
                   \ \Longleftrightarrow\ \br(F_E)=17.}
\]
The forward implication uses the audited exhaustive classification in
Theorem~\ref{thm:prior}; the reverse implication does not require another
classification.
\end{theorem}
\begin{proof}
If $\br(M_3)=17$, the simultaneous Borel span data have the first pair plane in
\eqref{eq:canonicalpairs} after cyclic rotation. A pair span remains a limit of
17 rank-one matrix spans, so $E\in\calG_{17}(Z)$.

For any $r$-dimensional matrix space, its inclusion tensor has border rank $r$
if and only if the space belongs to $\calG_r$ of the rank-one matrix variety.
Indeed, a convergent basis of rank-one matrix spans supplies the first-factor
coefficients of a degeneration of the inclusion tensor. Conversely, conciseness
in the $r$-dimensional first factor makes the image of the flattening continuous
on a neighborhood of the limiting tensor. This proves the second equivalence.

Finally, apply the linear map $D_E\to C$ given by $e_{ik}\mapsto z_{ki}$ and
$e_{ab}\mapsto0$ to \eqref{eq:FE}. It gives $M_3$. Border rank cannot increase
under a linear map, so $\br(F_E)=17$ implies $\br(M_3)\leq17$. The inherited
lower bound gives equality.
\end{proof}

\subsection{Repeated diagonal blocks}

Reorder rows of $A$ by the inner index $j$ first, then $i$, and columns of $B$ by
the inner index first, then $k$. An arbitrary matrix in $E$ becomes
\begin{equation}\label{eq:blockmatrix}
 \begin{pmatrix}H&P&R\\0&H&0\\0&0&H\end{pmatrix},
 \qquad H\in M_3(\C),
\end{equation}
where $P,R$ are supported on rows 2,3 and columns 1,2. The dimensions are
$9+4+4=17$.

\begin{lemma}\label{lem:rank3}
Every matrix in \eqref{eq:blockmatrix} has rank at least $3\rank H$.
\end{lemma}
\begin{proof}
If $q=\rank H$, select a nonsingular $q\times q$ submatrix of $H$. Select the
corresponding rows and columns in each of the three diagonal blocks. The
resulting $3q\times3q$ submatrix is upper block triangular with three copies of
the same nonsingular submatrix on its diagonal. Its determinant is nonzero.
\end{proof}

\begin{proposition}[Reduced pair section]\label{prop:pairsection}
Scheme-theoretically,
\[
 Z\cap\PP E=\Seg(\PP A'\times\PP B').
\]
In particular, its linear span has dimension eight.
\end{proposition}
\begin{proof}
Set-theoretically, Lemma~\ref{lem:rank3} forces $H=0$ in any rank-one element
of $E$. The resulting rank-one elements are exactly the small Segre variety.

At a nonzero rank-one matrix $a\otimes b$, the affine tangent space of the
rank-one variety is $a\otimes B+A\otimes b$. Every matrix in it has rank at most
two. Its intersection with $E$ therefore again has $H=0$. Its intersection
with $A'\otimes B'$ is precisely
$a\otimes B'+A'\otimes b$, the tangent space of the small Segre variety. Thus
every projective local tangent dimension is four, equal to the local dimension
of the set-theoretic intersection. Each local ring is regular. The section is
therefore reduced and has exactly the asserted scheme structure.
\end{proof}

\subsection{A cactus obstruction, not a border obstruction}

The cactus rank $\crk(T)$ is the least length of a zero-dimensional scheme in
the appropriate Segre variety whose linear span contains $T$. This differs
from a limit of spans: a flat point-scheme limit can lose span dimension.
The distinction is central in border apolarity.\src{Buczy\'nska--Buczy\'nski,
\emph{Apolarity, border rank and multigraded Hilbert scheme}}{5}

\begin{corollary}\label{cor:FEcactus}
The tensor $F_E$ satisfies $\crk(F_E)\geq18$.
\end{corollary}
\begin{proof}
Suppose a scheme of length at most 17 in
$\Seg(\PP D_E\times\PP A\times\PP B)$ spans $F_E$. Its scheme-theoretic image
in $Z$ has length at most 17 and its span contains the flattening image $E$.
A scheme of length at most 17 spans a vector space of dimension at most 17.
Since $\dim E=17$, this image would have to span exactly $E$ and hence be
contained in $Z\cap\PP E$. Proposition~\ref{prop:pairsection} places its span
in the eight-dimensional core, a contradiction.
\end{proof}

Corollary~\ref{cor:FEcactus} cannot be substituted for $\br(F_E)\geq18$.
Section~\ref{sec:control} gives a verified counterexample to that inference in
exactly the same tensor format.

\section{The full scheme section of the 17-plane}

Let $Y=\Seg(\PP A'\times\PP B'\times\PP C')$, so $\Span Y=\PP L$ and
$\dim Y=5$. Use coordinates $\lambda,z_1,\ldots,z_{16}$ on $\calS$, with
$\lambda$ the coefficient of $M_3$. Write $J_Y$ for the ideal of the small
Segre variety in $\C[z_1,\ldots,z_{16}]$.

\begin{theorem}[Restricted ideal and saturation]\label{thm:section}
The restriction of the homogeneous Segre ideal to $\PP\calS$ is
\begin{equation}\label{eq:restrictedideal}
 I_X|_{\calS}=J_Y+\lambda(\lambda,z_1,\ldots,z_{16}).
\end{equation}
Its saturation is $J_Y+(\lambda)$. Consequently
\[
 X\cap\PP\calS=Y
\]
scheme-theoretically, with $Y$ reduced.
\end{theorem}
\begin{proof}
The Segre ideal is generated by the $2\times2$ minors of its flattenings.
Restrict a tensor coordinate to $\calS$. It is either zero, $\lambda$, or one
of the 16 distinct variables $z_i$, because the two support sets are disjoint.
The restriction of any flattening minor, after setting $\lambda=0$, is a
small-Segre minor or zero. Thus every restricted minor belongs to the right
side of \eqref{eq:restrictedideal}.

Conversely, all small-Segre minors are restrictions of large-Segre minors.
Here are explicit certificates for the extra 17 monomials, in zero-based
row-major coordinate indices. The tensor coordinates of $M_3$ are
$(3i+j,3j+k,3k+i)$. In the $A$-flattening, use diagonal coordinates
$(0,0,0)$ and $(4,3,1)$. The off-diagonal product vanishes, so this minor is
$\lambda^2$. For every core coordinate $\beta$, use diagonal coordinates
$(0,0,0)$ and $\beta$. The off-diagonal product again vanishes, giving
$\lambda z_\beta$. These certificates prove equality.

The product of $\lambda$ with every degree-one coordinate belongs to the
ideal, so $\lambda$ lies in its saturation. The quotient after adding
$\lambda$ is the homogeneous coordinate ring of the reduced Segre variety
$Y$, whose ideal is saturated. This proves the saturation formula.
\end{proof}

The exact quadratic-space dimensions are
\[
 \dim (J_Y)_2=46,\qquad
 \dim (I_X|_{\calS})_2=63=46+17.
\]
The unsaturated coordinate ring has a one-dimensional extra degree-one class
annihilated by the homogeneous maximal ideal. Its diagonal Hilbert function
starts
\[
 1,\ 17,\ 90,\ 320,\ 875,\ 2016,\ 4116,\ 7680,\ 13365.
\]
For $d\geq2$ the value is $(d+1)^2\binom{d+3}{3}$, the Hilbert function of $Y$.
The extra class is at the vertex of the affine cone; it is not a nonreduced
projective component of $X\cap\PP\calS$.

\begin{corollary}\label{cor:schemelimit}
Any flat limit of the 17 point schemes accompanying a span limit to $\calS$
is contained scheme-theoretically in $Y$. No length-17 scheme in $X$ spans
$\calS$. The four-factor tensor $Q$ has cactus rank at least 18.
\end{corollary}
\begin{proof}
Containment of a scheme in its limiting linear span is closed, so a limiting
scheme is contained in $X\cap\PP\calS=Y$. Every scheme in $Y$ spans a subspace
of $L$, whose dimension is 16, not 17. For $Q$, project a hypothetical
length-17 spanning scheme to the last three factors and use first-factor
conciseness, exactly as in Corollary~\ref{cor:FEcactus}.
\end{proof}

\subsection{All positive prolongations}

For $a,b,c\geq1$, let $K_{abc}$ be the space of partially symmetric tensors of
multidegree $(a,b,c)$ all of whose contractions to degree $(1,1,1)$ belong to
$\calS$. Equivalently, it is the annihilator in that degree of the ideal
generated by $\calS^\perp$ in degree $(1,1,1)$.

\begin{proposition}\label{prop:prolong}
The first positive prolongations are
\begin{align*}
 K_{211}&=\Sym^2 A'\otimes B'\otimes C',&\dim K_{211}&=24,\\
 K_{121}&=A'\otimes\Sym^2 B'\otimes C',&\dim K_{121}&=40,\\
 K_{112}&=A'\otimes B'\otimes\Sym^2 C',&\dim K_{112}&=24.
\end{align*}
More generally, for $a+b+c\geq4$,
\[
 K_{abc}=\Sym^a A'\otimes\Sym^b B'\otimes\Sym^c C'.
\]
\end{proposition}
\begin{proof}
An element of $A\otimes\calS$ has a unique expression $u\otimes M_3+H$ with
$H\in A\otimes L$. Project the $B$ factor to $B/B'$. The second summand
disappears. If the two $A$ factors are symmetric, then
$u\otimes\overline M_3$ is symmetric in them. If $u\ne0$, this forces the
$A$-flattening image of $\overline M_3$ to lie in $\Span(u)$: contract the
other factors and use $u\otimes v=v\otimes u$ only when $v$ is proportional
to $u$. But that flattening has rank nine. Each of the nine slices has a
nonzero coordinate, and their supports are disjoint. Thus $u=0$.

Symmetry of $H\in A\otimes L$ then forces the extra $A$ factor into $A'$.
This proves the first formula. The other two follow by the same argument,
using a quotient in one of the other factors which preserves at least two
independent slices (in fact nine). The supplied wedge-kernel computations
check all three statements over $\Q$; their ranks are $129,113,129$ on
153-dimensional domains.

For higher degree, contract all but one additional slot. The degree-four
statements show that every resulting coefficient is supported on the small
factor spaces. Repeating this for each slot puts every factor in the
corresponding small space. The reverse inclusion is immediate.
\end{proof}

These formulas describe necessary prolongation spaces. They do not force a
limiting 17-dimensional degree-$(1,1,1)$ span to equal the span of the
saturated limiting point scheme.

\section{Exact rank-17 certificates after grouping factors}

The next construction is a positive result for \emph{grouped} tensors. It
also gives an explicit control against an invalid cactus-to-border inference.

\subsection{A Fourier compression formula}

Let $p,q,s$ be positive integers, $N=pq+1$, and $p\leq s\leq N-q$.
Consider independent basis vectors $d_0,d_{rj}$ in an $N$-dimensional space,
$e_k,e'_r$ in an $(s+q)$-dimensional space, and $c_k$ in an $s$-dimensional
space. Define
\begin{equation}\label{eq:compression}
 T_{p,q,s}=d_0\otimes\sum_{k=1}^{s}e_k\otimes c_k
       +\sum_{r=1}^{q}\sum_{j=1}^{p}d_{rj}\otimes e'_r\otimes c_j.
\end{equation}

\begin{theorem}\label{thm:compression}
The tensor in \eqref{eq:compression} has border rank exactly $N$.
\end{theorem}
\begin{proof}
Work with residue labels modulo $N$. Set
\[
 J=\{0,\ldots,p-1\},\qquad R=\{p,2p,\ldots,qp\}.
\]
Choose $K$ of size $s$ containing $J$ and disjoint from $R$, and put $H=K\cup R$.
Define $w_j=0$ on $J$, $w_h=2$ on $R$, and $w_h=1$ on $K\setminus J$.
Use the residue labels to index the corresponding $e$ and $c$ basis vectors.
The differences $r-j$ for $r\in R,j\in J$ are exactly the distinct integers
$1,\ldots,N-1$, so relabel $d_{rj}$ by $d_{r-j}$.

For a primitive $N$th root of unity $\zeta$, put
\begin{equation}\label{eq:fourier}
 \frac1N\sum_{a=0}^{N-1}
 \left(d_0+t^2\sum_{\ell=1}^{N-1}\zeta^{a\ell}d_\ell\right)
 \otimes\left(\sum_{h\in H}t^{-w_h}\zeta^{-ah}e_h\right)
 \otimes\left(\sum_{k\in K}t^{w_k}\zeta^{ak}c_k\right).
\end{equation}
The root average is zero unless $\ell-h+k=0$ modulo $N$, when it is one.
For $\ell=0$, only $h=k$ occurs, with power zero, giving the first summand
of \eqref{eq:compression}. For $\ell\ne0$, the power is
$2+w_k-w_h\geq0$. It is zero exactly when $h\in R,k\in J$, giving all the
second summands. All remaining terms have positive powers. Thus
\eqref{eq:fourier} is an $N$-term degeneration of \eqref{eq:compression}.

The flattening in the $d$ factor has rank $N$: the core matrix units are
independent, and the diagonal $e_k\otimes c_k$ sum has disjoint support.
This gives the matching lower bound.
\end{proof}

The verifier works at $N=17$ and checks root orthogonality in
$\Q[\zeta]/(1+\zeta+\cdots+\zeta^{16})$. It expands the complete coefficient
tensor, not just its prescribed nonzero entries.

\subsection{Application to the four-factor target}

In $Q$, keep $D$ separate and keep one of $A,B,C$ as a singleton. Group the
other two matrix factors. The nine slices of $M_3$ in the grouped factor are
independent and disjoint in support from its $q$ core coordinate vectors.
Thus the grouped tensor is linearly equivalent to \eqref{eq:compression}:
\begin{center}
\begin{tabular}{lrrrrr}
\toprule
Singleton&$p$&$q$&$s$&Grouped factor dimension&$N$\\
\midrule
$A$&2&8&9&17&17\\
$B$&4&4&9&13&17\\
$C$&2&8&9&17&17\\
\bottomrule
\end{tabular}
\end{center}

\begin{corollary}\label{cor:grouped}
The three grouped border ranks of $Q$ just described are exactly 17.
\end{corollary}

The coefficient expansions have residual degrees $1,2,3$. Their total
coordinate $\ell^1$ residual coefficients are $272,236,272$, respectively,
so the same numbers times $|t|$ are Frobenius error bounds for $0<|t|\leq1$.
At $t=1$ and Fourier index zero, each grouped middle factor, reshaped in the
two original factors, has matrix rank four. These particular middle factors
are not decomposable. The formula is therefore not a four-factor
border-rank-17 decomposition of $Q$.

\subsection{A control tensor with border rank below cactus rank}\label{sec:control}

Let $a_0,\ldots,a_8$ and $b_0,\ldots,b_8$ be bases and set
\begin{equation}\label{eq:control}
 W=d_0\otimes\sum_{i=0}^{8}a_i\otimes b_i
  +\sum_{i=0}^{7}\sum_{j=0}^{1}d_{1+2i+j}\otimes a_i\otimes b_j
       \in\C^{17}\otimes\C^9\otimes\C^9.
\end{equation}
This is a different tensor from $F_E$.

\begin{theorem}[Exact counterexample to a tempting inference]\label{thm:wild}
The tensor $W$ has $\br(W)=17$ and $\crk(W)\geq18$.
\end{theorem}
\begin{proof}
Apply Theorem~\ref{thm:compression} at $(p,q,s)=(2,8,9)$, sending the nine
$e_k$ to $a_0,\ldots,a_8$ and the eight $e'_r$ to $a_0,\ldots,a_7$. This
linear restriction gives \eqref{eq:control}, so its border rank is at most 17.
The $D$-flattening has rank 17, giving equality. The supplied expansion has
25 constant unit coordinates and residual coordinate norms $70,44,14$ in
degrees $1,2,3$.

The matrix plane of $W$ is
$E_W=\Span(I_9)\oplus(\Span(a_0,\ldots,a_7)\otimes\Span(b_0,b_1))$.
A matrix with nonzero $I_9$ coefficient has rank at least seven: its columns
2 through 8 are independent scalar multiples of the corresponding standard
basis vectors. Hence its rank-one locus is the core Segre variety, whose
span has dimension 16. The same rank-at-least-seven versus tangent-rank-at-
most-two argument as in Proposition~\ref{prop:pairsection} proves that the
section is reduced. First-factor conciseness then excludes a spanning
scheme of length 17, as before.
\end{proof}

This control is included to prevent a false completion of AC-03 from
Corollary~\ref{cor:FEcactus}. It is not used to assert that $F_E$ has border
rank 17. No claim of publication-level novelty is made for this compression
family or control; related constructive minimal-border-rank and wild-tensor
phenomena occur in the recent literature.\src{Kassabov--Landsberg--Souza--Speegle,
\emph{Disjoint and nearly disjoint sums of matrix multiplication tensors and
their centroids}}{7}

\section{Necessary conditions for cubic-order constructions}

Suppose 17 normalized holomorphic factor curves have limiting points in $Y$:
\[
 a_s(t)\otimes b_s(t)\otimes c_s(t),\qquad 1\leq s\leq17.
\]
Normalization means that each factor curve has a nonzero limit. Scalar
coefficients are allowed poles. Their order is measured \emph{after} this
normalization; shifting powers between a factor and its scalar would
otherwise make a pole-order statement meaningless.

Let $\pi$ project all three factors to $A/A',B/B',C/C'$. Then
\begin{equation}\label{eq:projection}
\begin{split}
\pi(M_3)={}&\bar x_{11}\otimes\bar y_{11}\otimes\bar z_{11}
 +\bar x_{11}\otimes\bar y_{12}\otimes\bar z_{21}\\
&+\bar x_{32}\otimes\bar y_{23}\otimes\bar z_{33}
 +\bar x_{33}\otimes\bar y_{33}\otimes\bar z_{33}.
\end{split}
\end{equation}
Its three flattening ranks are $(3,4,3)$.

\begin{lemma}\label{lem:pole2}
No such expression for $M_3$ can have all scalar poles of order at most two.
\end{lemma}
\begin{proof}
Each projected factor is $O(t)$, so every projected decomposable term is
$O(t^3)$. Poles of order at most two leave a sum tending to zero, contrary
to \eqref{eq:projection}.
\end{proof}

For the cubic test, assume the 17 base points have 16 independent triple
products with a unique relation whose 17 coefficients are nonzero. Identify
functions on the 17 points with the semisimple algebra $R=\C^{17}$. Let
$A_0,B_0,C_0\subset R$ be the factor-coordinate function spaces, of dimensions
$2,4,2$, and write the relation as a functional $\lambda:R\to\C$ with
\[
 \lambda(A_0B_0C_0)=0.
\]
The bilinear form $(f,g)\mapsto\lambda(fg)$ is nondegenerate because all
relation coefficients are nonzero.

The first normal function spaces are
\begin{equation}\label{eq:normalspaces}
 U_A=\ann(B_0C_0)/A_0,\quad
 U_B=\ann(A_0C_0)/B_0,\quad
 U_C=\ann(A_0B_0)/C_0,
\end{equation}
where annihilators use this bilinear form. When the pair product dimensions
are $8,4,8$, these spaces have dimensions $7,9,7$.

\begin{lemma}[Second-normal conditions]\label{lem:normal}
A cubic-pole expression for $M_3$ at these base points supplies subspaces of
$U_A,U_B,U_C$ of dimensions at least $(3,4,3)$ satisfying
\begin{equation}\label{eq:pairnormal}
 \lambda(\mathsf A\mathsf B C_0)=0,\qquad
 \lambda(\mathsf A B_0\mathsf C)=0,\qquad
 \lambda(A_0\mathsf B\mathsf C)=0.
\end{equation}
The all-normal constant tensor is the restriction of
$(f,g,h)\mapsto\lambda(fgh)$ to the first normal images. It must equal
\eqref{eq:projection} after the physical coordinate maps.
\end{lemma}
\begin{proof}
The coefficient of $t^{-3}$ is a relation among the base products, hence a
scalar multiple of the unique $\lambda$. If it were zero, Lemma~\ref{lem:pole2}
would apply. Project one factor outside its small space. The coefficient of
$t^{-2}$ then says that each first normal coordinate annihilates the product
of the other two base coordinate spaces. This gives \eqref{eq:normalspaces}.

Project two factors outside. The coefficient of $t^{-1}$ involves exactly
their first normal coordinates and the base value of the third factor,
which gives \eqref{eq:pairnormal}. Project all three factors: the constant
term contains exactly three first normal coordinates and the leading
scalar relation. Its flattening ranks must be those of \eqref{eq:projection}.

Passing to the quotient representatives in \eqref{eq:normalspaces} is
legitimate. Terms added from a base space vanish in the second-normal
conditions by the first conditions; they vanish in the cubic tensor by
\eqref{eq:pairnormal}. Thus all the stated restrictions are well defined.
\end{proof}

\subsection{An exterior-contraction certificate}

For a bilinear form $m$ on $U\times V$, define
\[
 \partial_m(u_1\wedge\cdots\wedge u_a\otimes
             v_1\wedge\cdots\wedge v_b)
 =\sum_{i,j}(-1)^{i+j}m(u_i,v_j)
     u_1\wedge\cdots\widehat u_i\cdots\wedge u_a
       \otimes v_1\wedge\cdots\widehat v_j\cdots\wedge v_b.
\]
Numbering signs from zero instead of one gives the same formula. Apply all
the forms in \eqref{eq:pairnormal} to
\begin{equation}\label{eq:exteriordomain}
 \mathcal V=\Lambda^3U_A\otimes\Lambda^4U_B\otimes\Lambda^3U_C,
 \qquad\dim\mathcal V=\binom73\binom94\binom73=154\,350.
\end{equation}
Let $\mathcal D$ be their joint contraction map.

\begin{lemma}\label{lem:exterior}
If $\mathcal D$ is injective, no subspaces of dimensions $(3,4,3)$ satisfy
\eqref{eq:pairnormal}. Therefore no cubic-pole expression for $M_3$ exists
at the specified base configuration.
\end{lemma}
\begin{proof}
Bases of such subspaces give a nonzero decomposable product of three
Pl\"ucker vectors in \eqref{eq:exteriordomain}. Every contraction is zero by
\eqref{eq:pairnormal}, contradicting injectivity.
\end{proof}

\section{All normalized cyclic-root base configurations}

Take $R=\C[z]/(z^{17}-1)$ and choose exponent sets $A_0,B_0,C_0$ of sizes
$2,4,2$ whose triple sums are 16 distinct residues. The 17 evaluation points
at the roots of unity give 16 independent base products. If $m$ is the
missing residue, the unique relation is $\lambda(f)=[z^m]f$. Its evaluation
weights are nonzero. Products and pairings have rational structure constants
in the residue basis.

Translations of the three exponent sets, interchange of their basis labels,
and rescaling $z\mapsto z^a$, $a\ne0\pmod {17}$, normalize the sets to
\begin{equation}\label{eq:rootnorm}
 A_0=\{0,1\},\quad B_0=\{0,b_1,b_2,b_3\},\quad C_0=\{0,c\},
 \quad 1\leq b_1<b_2<b_3\leq16.
\end{equation}
Exactly 24 normalized configurations have 16 distinct triple sums. This is
a finite enumeration of \emph{exponent configurations}, not an enumeration
of arbitrary complex point configurations.

\begin{theorem}[Cyclic cubic obstruction]\label{thm:cyclic}
For every configuration in \eqref{eq:rootnorm} with 16 distinct triple sums,
the map $\mathcal D$ in \eqref{eq:exteriordomain} has full column rank
154\,350 over $\Q$ and $\C$. Hence no normalized holomorphic factor curves
with those limiting points and scalar poles of order at most three can
express $M_3$ in the limit.
\end{theorem}
\begin{proof}[Computer-assisted proof]
For each configuration, the normal quotient spaces are obtained exactly
from the rule $\lambda(z^az^b)=1$ when $a+b=m$ modulo 17 and zero otherwise.
Their dimensions are $7,9,7$. The coefficient of a contraction column is
$0,1$, or $-1$.

First apply the two $AB$ forms to
$\Lambda^3U_A\otimes\Lambda^4U_B$, of dimension 4410. The computation splits
by the sum of character labels modulo 17. It constructs a kernel basis
modulo 101 and independently checks every kernel relation against the
original columns. Tensor that kernel with $\Lambda^3U_C$, then apply the
remaining four $AC$ and two $BC$ forms. Each residual character block has
full column rank. The combined rank is exactly 154\,350 in each of the 24
cases.

A full-column-rank reduction modulo a prime certifies a nonzero maximal
minor of the original integer matrix. It therefore certifies injectivity
in characteristic zero. There is no claim that testing complex subspaces
can be replaced by enumerating subspaces over a finite field.
Lemma~\ref{lem:exterior} completes the proof.
\end{proof}

The two configurations
\[
 (\{0,1\},\{0,2,4,6\},\{0,8\}),\qquad
 (\{0,1\},\{0,2,8,10\},\{0,4\})
\]
are independently checked modulo 103 as well. Appendix~\ref{app:hand}
also gives a hand proof for the first, independent of the exterior rank
calculation.

\begin{corollary}[A generic reduced-configuration obstruction]\label{cor:generic}
There is a nonempty Zariski-open subset of the ordered 17-point configurations
in $Y$ for which cubic-pole expressions for $M_3$ are impossible.
\end{corollary}
\begin{proof}
Restrict to the nonempty open set where triple evaluation has rank 16,
the unique relation has no zero coefficient, and the pair product spaces
have maximal dimensions. The normal spaces in \eqref{eq:normalspaces} form
vector bundles of ranks $7,9,7$. On local charts their contractions depend
algebraically on the points. Injectivity of $\mathcal D$ is open. The first
configuration in Theorem~\ref{thm:cyclic} proves this open set is nonempty.
The configuration variety is irreducible, so the open subset is also dense.
\end{proof}

The corollary does not cover the exceptional closed locus or higher poles.
In particular, the grid in the next section has a nonzero exterior kernel,
so a blanket injectivity statement would be false.

\section{The grid-plus-barycenter obstruction}

Take the 16 coordinate product points of $Y$ and the product of the three
barycenter vectors. Normalize each barycenter factor so its coordinates
sum to one. The product at the barycenter is $1/16$ times the sum of the 16
grid products, so the unique relation has weights $-1$ on the grid and
$16$ on the barycenter.

Use the group $G=(\mathbb Z/2)^4$ to label functions on the grid. The constant
function is one on all 17 points; each nontrivial character is extended by
zero at the barycenter. A Fourier change of basis within the factors gives
base character spaces
\begin{equation}\label{eq:gridbase}
 A_0=\langle1,\chi_1\rangle,\quad
 B_0=\langle1,\chi_2,\chi_4,\chi_6\rangle,\quad
 C_0=\langle1,\chi_8\rangle.
\end{equation}
For nontrivial characters, a relation moment is $-16$ when the binary
exclusive-or sum of the labels is zero and otherwise zero. The irrelevant
common factor $-16$ is dropped in the matrices.

The three normal quotient bases are
\begin{align*}
 U_A&=(3,5,7,9,11,13,15),\\
 U_B&=(3,5,7,10,11,12,13,14,15),\\
 U_C&=(9,10,11,12,13,14,15).
\end{align*}
The joint exterior map is \emph{not} injective here. Its kernel has dimension
135, and genuine $(3,4,3)$ subspaces satisfying the pair conditions exist.
Appendix~\ref{app:gridcontrol} displays one. Further equations are necessary.

\begin{theorem}[Grid cubic obstruction]\label{thm:grid}
For the configuration \eqref{eq:gridbase}, every triple of normal subspaces
of dimensions $(3,4,3)$ satisfying \eqref{eq:pairnormal} has identically zero
cubic form $\lambda(fgh)$ on their product. Consequently the grid-plus-
barycenter configuration cannot be the base of a cubic-pole expression
for $M_3$.
\end{theorem}
\begin{proof}[Exact certificate and its algebraic interpretation]
Let $\mathcal K=\ker\mathcal D\subset\mathcal V$.
The proof uses five finite steps; all polynomial operations after the
initial modular rank bound are over $\Q$.

\textbf{The complete linear kernel.}
The file \texttt{grid\_joint\_kernel.json} gives 135 integral vectors in
$\mathcal V$. Each is substituted into the original contraction columns,
and its image is checked to be zero over $\Q$. Their independence is checked
by rational elimination. Independently, a block computation modulo 101
gives rank $154\,215=154\,350-135$ for the original integer map. A
characteristic-zero kernel has dimension at most the modular kernel
dimension. The 135 independent rational vectors therefore form the entire
kernel over $\Q$ and, by scalar extension, over $\C$.

\textbf{Necessary quadratic equations.}
A candidate product of Pl\"ucker vectors must be decomposable across the three
exterior factors and must have a Grassmannian Pl\"ucker vector in each factor.
Thus its coordinates satisfy the three-factor Segre minors and the
Grassmann Pl\"ucker relations on every fixed fiber. Substitute a general
linear combination of the 135 kernel vectors into these equations.
Two batches of explicit Segre minors reduce to nonzero scalar multiples of
individual parameter squares and eliminate 87 parameters, leaving 48.
The selected, explicitly recorded necessary quadrics span a
1143-dimensional subspace of the 1176-dimensional space of quadrics in
those 48 parameters. It is not necessary to assert that these selected
quadrics generate the whole ideal.

\textbf{Linear consequences of quadratic ideal membership.}
There are 17 independent linear forms $\ell$ whose products $\ell x_j$ with
every remaining parameter lie in the selected quadratic span. The verifier
computes them as a rational kernel and checks those memberships. Therefore
$\ell^2$ belongs to that span, and every complex solution has $\ell=0$.
Restricting to their common zero space leaves 31 parameters. The restricted
quadratic span has rank 463.

\textbf{The cubic contraction.}
Contract one exterior slot in each factor against
$\tau(f,g,h)=\lambda(fgh)$. This defines
\[
 \partial_\tau:\mathcal V\longrightarrow
 \Lambda^2U_A\otimes\Lambda^3U_B\otimes\Lambda^2U_C.
\]
On the remaining 31-parameter linear space its coefficient functionals
span a seven-dimensional space. Three independent such functionals have
squares in the restricted necessary quadratic span; hence they vanish at
every complex solution. After imposing them, 28 parameters remain and the
image of $\partial_\tau$ lies in a four-dimensional tensor space.

\textbf{An exact flattening contradiction.}
For a basis $V_1,\ldots,V_4$ of that four-dimensional image, each $V_i$ has
23 columns in its first-factor flattening which occur in no other $V_j$.
Those exclusive columns have rank nine. Therefore every nonzero linear
combination of the $V_i$ has first flattening rank at least nine: choose a
nonzero coefficient and restrict to its exclusive columns.

On the other hand, if the original exterior vector is a product of
Pl\"ucker vectors of spaces $\mathsf A,\mathsf B,\mathsf C$ of dimensions
$3,4,3$, its cubic contraction lies in
\[
 \Lambda^2\mathsf A\otimes\Lambda^3\mathsf B\otimes\Lambda^2\mathsf C.
\]
Its first flattening rank is at most $\dim\Lambda^2\mathsf A=3$.
It must therefore be zero. The cofactor wedge bases in these three spaces
are independent, so this is equivalent to $\tau$ being zero on
$\mathsf A\times\mathsf B\times\mathsf C$.

If larger first normal image spaces existed for $M_3$, any nonzero cubic
coordinate could be extended to subspaces of dimensions $(3,4,3)$, since
the projected flattening ranks are $(3,4,3)$. The preceding conclusion
would make that coordinate zero. Thus Lemma~\ref{lem:normal} contradicts
\eqref{eq:projection} in every case.
\end{proof}

This proof neither infers complex nonexistence from unsuccessful numerical
searches nor treats a modular kernel as an exact rational kernel without
checking it. The integral basis plus the one-sided modular rank bound
provides the required characteristic-zero equality. The additional
quadratic equations are essential: the pair conditions alone admit the
control in Appendix~\ref{app:gridcontrol}.

\section{A separate restricted monomial Fourier exclusion}

There is a different root-of-unity search for the pair plane $E$, distinct
from the cubic-base tests above. Consider curves indexed by $s\in\mathbb Z/17$
whose every matrix coordinate has one nonzero character monomial:
\[
 a_{ij}^{(s)}(t)=c_{ij}t^{\alpha_{ij}}\zeta^{s e_{ij}},\qquad
 b_{\ell k}^{(s)}(t)=d_{\ell k}t^{\beta_{\ell k}}\zeta^{s f_{\ell k}},
\]
with all $c_{ij},d_{\ell k}\ne0$. Arbitrary Laurent or rational exponents may
be allowed; the obstruction below already holds for real exponents.

A span limit equal to $E$ forces the nine repeated diagonal slices and the
eight core units to occupy 17 distinct characters. Equality of characters
within a diagonal slice implies
\[
 e_{ij}=u_i+v_j,\qquad f_{jk}=w_k-v_j\pmod {17}.
\]
Equality of valuations within each such slice likewise gives
$\alpha_{ij}=a_i+b_j$ and $\beta_{jk}=c_k-b_j$.
Normalize $u_0=w_0=v_0=0$, $u_1=1$, and interchange the last two inner labels
so $v_1<v_2$. There are 810 possible distinct nine-character main sumsets;
exactly 24 character patterns also assign eight disjoint characters to the
core units.

Every off-target matrix coordinate in a given character must have strictly
larger valuation than that character's prescribed limiting slice or unit.
There are 46 such strict inequalities in six independent real weights.

\begin{theorem}[Restricted monomial obstruction]\label{thm:toric}
No curves of the preceding single-character, single-monomial form have
17-dimensional pair spans tending to $E$.
\end{theorem}
\begin{proof}[Exact finite certificate]
The verifier enumerates all 810 main sumsets and all 24 normalized full
patterns. For each full pattern it constructs the 46 integer gap rows
$g_i\in\mathbb Z^6$. The certificate supplies positive rational coefficients
$q_i$, on a nonempty subset of these rows, satisfying
\[
 \sum_iq_i=1,\qquad \sum_iq_i g_i=0.
\]
If all required gaps $g_i\cdot w$ were positive, their positive weighted
sum would be positive, contradicting the displayed zero identity. Every
row and rational dependence is regenerated and checked. Thus all 24
patterns are impossible, even over arbitrary real valuation weights.
\end{proof}

Nonzero amplitudes cannot cancel an off-target coordinate in a fixed
character column: it is a single product of two nonzero amplitudes.
This justifies the strict-gap necessity. The theorem does not cover sums
of character monomials in a coordinate, arbitrary factor curves, or
non-diagonal changes of the ambient coordinates during the degeneration.
The coincidence that this enumeration and Section 6 each contain 24
normalized patterns does not make them the same classification.

\section{What remains to determine the exact border rank}

The new pair reduction sharpens the unresolved rank-17 question to
\begin{equation}\label{eq:remaining}
 E\stackrel{?}{\in}\calG_{17}(\Seg(\PP A\times\PP B)),
 \qquad\text{equivalently}\qquad \br(F_E)\stackrel{?}{=}17.
\end{equation}
A positive certificate must provide a genuine limiting span or an exact
17-term degeneration of $F_E$. A universal nonmembership proof would
establish $\br(M_3)\geq18$. It would \emph{not} by itself determine whether
the answer is 18, 19, or 20.

The results here eliminate a monomial Fourier route and cubic constructions
at the specified cyclic and grid bases. They do not eliminate arbitrary
special reduced configurations, nonreduced limits, or higher normalized
scalar poles. Corollaries about cactus rank do not fill that gap, as the
control tensor proves. The positive grouped certificates do not fill it
either, because grouping forgets the required decomposability of a merged
factor.

A separate exploratory route considers unrestrictions to concise
minimal-border-rank tensors and necessary commutative-algebra conditions.
The general framework is given by Jagie\l\l a--Jelisiejew.\src{Jagie\l\l a--Jelisiejew,
\emph{Unrestrictions and concise secant varieties}, version 2}{6}
The exploratory basis-line enumeration in this package is incomplete at
parameter-bearing branches and is not used for an exclusion. Finite-field
samples and floating-point optimization records are likewise not proofs of
nonexistence over $\C$.

The exact numerical interval is unchanged. None of the arguments in this
report selects a value from $17,18,19,20$, and the package must not be
represented as a completed solution of AC-03.

\section{Reproducibility and audit boundaries}

From the extracted directory, run
\begin{Verbatim}[fontsize=\small]
python code/check_manifest.py
python -m pip install -r requirements.txt
python code/verify_all.py
\end{Verbatim}
The new certificates use only Python's standard library. SymPy 1.14.0 is
needed for the retained prior verifier. The tested interpreter is Python
3.13.5. Do not use Python's \texttt{-O} option: assertions are verification
checks, and the combined runner rejects optimized mode.

The combined runner regenerates the scheme-section, pair-plane, grouped,
control, cyclic, monomial, exterior, and grid results. It runs eight
consistency tests and independently checks the two principal cyclic
configurations modulo 103. It extracts the unchanged previous archive into
a temporary directory, verifies its manifest, and recomputes its complete
round-two and round-one checks. The optional \texttt{--skip-prior} switch
checks only the new work and must not be described as revalidating the
inherited bounds or classification.

\begin{center}
\begin{tabular}{p{0.37\linewidth}p{0.55\linewidth}}
\toprule
File&Role\\
\midrule
\texttt{code/pair\_geometry.py}&35-coordinate target, restriction, block conventions.\\
\texttt{code/section\_geometry.py}&Restricted ideal and complete first prolongation kernels.\\
\texttt{code/grouped\_degeneration.py}&All three exact grouped length-17 expansions.\\
\texttt{code/wild\_control.py}&Independent border-17/cactus-at-least-18 control.\\
\texttt{code/exterior\_obstruction.py}&Complete cyclic-root normalization and integer-rank certificates.\\
\texttt{code/grid\_obstruction.py}&Integral kernel, necessary quadrics, and cubic flattening obstruction.\\
\texttt{code/toric17\_verify.py}&All monomial patterns and positive rational gap certificates.\\
\texttt{certificates/}&Exact input witnesses, not expected numerical answers alone.\\
\texttt{results/}&Regenerated proof records and verification logs.\\
\texttt{research/}&Exploratory calculations, with incompleteness and sampling boundaries.\\
\bottomrule
\end{tabular}
\end{center}

The report contains the logical reductions interpreting the finite
calculations. Universal rank and scheme arguments are not inferred from
sample matrices. The modular computations prove nonzero integer minors or
bound a rational kernel which is independently constructed; they do not
enumerate complex solutions. The checks are computer-assisted mathematics,
not proof-assistant formalization, independent peer review, or a guarantee
against every implementation error.

\appendix
\section{A hand proof for one cyclic configuration}\label{app:hand}

This appendix proves the cubic obstruction for
\[
 A_0=\langle1,z\rangle,\quad
 B_0=\langle1,z^2,z^4,z^6\rangle,\quad
 C_0=\langle1,z^8\rangle
\]
in $R=\C[z]/(z^{17}-1)$, with $\lambda=[z^{16}]$, without the large exterior
rank certificate. The quotient representatives are
\[
 U_A=z^3\langle1,z^2,\ldots,z^{12}\rangle,\quad
 U_B=\langle z,z^3,z^5,z^9,z^{10},z^{11},z^{12},z^{13},z^{14}\rangle,
 \quad U_C=\langle z,\ldots,z^7\rangle.
\]
Suppose spaces $\mathsf A,\mathsf B,\mathsf C$ of dimensions $3,4,3$ satisfy
the second-normal conditions.

Put $u=z^2$ and write $\mathsf A=z^3\mathcal A(u)$ with
$\mathcal A\subset\C[u]_{\leq6}$ of dimension three. The odd part of
$\mathsf C$ has the form $z\mathcal D(u)$ with
$\mathcal D\subset\C[u]_{\leq3}$. The even part is contained in
$\langle z^2,z^4,z^6\rangle$. The $AC$ conditions say that the coefficients
of $u^3,u^4,u^5,u^6$ in every product $ad$, $a\in\mathcal A,d\in\mathcal D$,
vanish. The even part contributes zero to these conditions by parity.

Pair $\C[u]_{\leq6}$ with itself by the coefficient of $u^6$. The condition
is $\mathcal A\perp\mathcal D\C[u]_{\leq3}$. If $e=\dim\mathcal D>0$, then
\[
 \dim(\mathcal D\C[u]_{\leq3})\geq e+3.
\]
For example, choose a basis of $\mathcal D$ with distinct leading degrees;
multiplication by $1,u,u^2,u^3$ produces at least $e+3$ distinct leading
positions. Since $\dim\mathcal A=3$ in a seven-dimensional perfect pairing,
$e+3\leq4$, so $e\leq1$.

If $e=0$, the three-dimensional $\mathsf C$ is the entire even subspace.
The $BC$ conditions with $A_0=\langle1,z\rangle$ give six independent
conditions on the nine-dimensional $U_B$. Thus $\dim\mathsf B\leq3$, a
contradiction.

If $e=1$, choose $0\ne d\in\C[u]_{\leq3}$. The four forms imposing the
central coefficient conditions are independent: multiplication by $d$ on
$\C[u]_{\leq3}$ is injective and the coefficient pairing is perfect. Hence
\[
 \mathcal A=\mathcal A_d:=\{a\in\C[u]_{\leq6}:
                       [u^3]ad=\cdots=[u^6]ad=0\},\qquad\dim\mathcal A_d=3.
\]
We claim the six-dimensional sum
$z^3\mathcal A_d+z^{11}\mathcal A_d$ is direct and disjoint from
$P=\langle1,z,z^8,z^9\rangle$.

Indeed, suppose $z^3a(z^2)+z^{11}b(z^2)\in P$ modulo $z^{17}-1$.
Comparison of the even and odd residues forces
\[
 b=b_{\mathrm{low}}+\beta u^3,\quad \deg b_{\mathrm{low}}\leq2,
 \qquad a=\alpha u^3-u^4b_{\mathrm{low}}.
\]
Since $b\in\mathcal A_d$ and $\deg(bd)\leq6$, its central coefficients
vanish only if $\deg(bd)\leq2$. The polynomial ring is a domain, so $\beta=0$.
Now
\[
 ad=u^3(\alpha d-u b_{\mathrm{low}}d).
\]
The polynomial in parentheses has degree at most three. Its coefficients
must all vanish, so $d(\alpha-u b_{\mathrm{low}})=0$. Again using the domain
property gives $\alpha=0$ and $b_{\mathrm{low}}=0$, and then $a=b=0$.
This proves the claim.

Consequently
\[
 \dim\big(A_0C_0+\mathsf A C_0\big)=4+6=10.
\]
The first and second normal conditions put the full $B$ function space in
its annihilator, of dimension $17-10=7$. This annihilator contains $B_0$,
which has dimension four. Its quotient therefore has dimension at most
three, again contradicting $\dim\mathsf B=4$.

This proof covers every nonzero complex $d$. The 40 rational coefficient
examples checked by \texttt{cyclic\_third\_order.py} are implementation
controls, not a substitute for this universal argument.

\section{A grid control for the pair conditions}\label{app:gridcontrol}

In the ordered normal character bases of Section 7, the following columns
span spaces of dimensions $(3,4,3)$:
\[
 A_1=\begin{pmatrix}
 1&0&0\\1&0&0\\0&1&0\\0&0&1\\-1&0&0\\-1&0&0\\0&0&1
 \end{pmatrix},\qquad
 C_1=\begin{pmatrix}
 0&0&-1\\1&0&0\\-1&0&0\\-1&0&0\\1&0&0\\0&1&0\\0&0&1
 \end{pmatrix},
\]
\[
 B_1=\begin{pmatrix}
 1&0&0&0\\0&1&0&0\\0&0&0&0\\0&0&1&0\\
 \frac12&\frac12&\frac12&-\frac12\\0&0&0&1\\
 \frac12&\frac12&-\frac12&\frac12\\0&0&0&0\\0&0&0&0
 \end{pmatrix}.
\]
Direct rational substitution gives all eight families of pairwise
conditions in \eqref{eq:pairnormal}. For every triple of columns the cubic
moment is zero. The consistency suite checks both facts. This demonstrates
why merely extending the cyclic injectivity claim to the grid would be
incorrect; the additional quadratic and flattening argument is needed.

\clearpage
\section*{Sources}
\addcontentsline{toc}{section}{Sources}
\begingroup
\raggedright
\begin{enumerate}[label={[\arabic*]},leftmargin=2em]
\item Andrew T., \emph{OpenProblemsInNLA}, AC-03 statement, public repository.
Retrieved 14 September 2026.
\url{https://raw.githubusercontent.com/ajt60gaibb/OpenProblemsInNLA/main/arithmetic-and-complexity/AC-03/README.md}.
Used for the problem specification and published-status context.

\item Austin Conner, Alicia Harper, J. M. Landsberg,
\emph{New lower bounds for matrix multiplication and the $3\times3$ determinant},
Forum of Mathematics, Pi 11 (2023), e17. Preprint arXiv:1911.07981.
\url{https://doi.org/10.1017/fmp.2023.14}.
Source of the inherited lower bound and border-apolarity classification framework.

\item A. V. Smirnov, \emph{The bilinear complexity and practical algorithms
for matrix multiplication}, Computational Mathematics and Mathematical Physics
53(12) (2013), 1781--1795; Russian edition 53(12), 1970--1984.
\url{https://doi.org/10.1134/S0965542513120129}.
Table 6 of the Russian edition, pp. 1982--1983, is the source of the retained
20-term degeneration. Bibliographic record:
\url{https://www.mathnet.ru/eng/zvmmf9955}.

\item Tomasz Ma\'ndziuk, \emph{Border rank lower bounds beyond weak border
apolarity}, arXiv:2609.12121, submitted 10 September 2026.
\url{https://arxiv.org/abs/2609.12121}.
Source context for the square test in the preceding package, not a claimed
new numerical bound for AC-03.

\item Weronika Buczy\'nska and Jaros\l aw Buczy\'nski,
\emph{Apolarity, border rank and multigraded Hilbert scheme}, arXiv:1910.01944,
accepted version dated 7 November 2020.
\url{https://arxiv.org/abs/1910.01944}.
Used for the distinction between actual border-rank data and weaker apolar
or saturated-scheme conditions.

\item Jakub Jagie\l\l a and Joachim Jelisiejew,
\emph{Unrestrictions and concise secant varieties}, arXiv:2604.24879v2,
29 June 2026.
\url{https://arxiv.org/abs/2604.24879}.
Framework for the exploratory concise-extension route; it is not invoked
to claim that the required extensions have been classified here.

\item Martin Kassabov, J. M. Landsberg, Victor Souza, Philip Speegle,
\emph{Disjoint and nearly disjoint sums of matrix multiplication tensors and
their centroids}, arXiv:2608.27434, 27 August 2026.
\url{https://arxiv.org/abs/2608.27434}.
Related constructive and wild-tensor context. The displayed formulas in
this report are independently expanded by the supplied code.

\item Josh Alman and Baitian Li, \emph{Asymptotic Rank Speedup Theorems,
Revisited}, CCC 2026, LIPIcs 383, pp. 36:1--36:42.
\url{https://doi.org/10.4230/LIPIcs.CCC.2026.36}.
Asymptotic rank is distinct from the finite-size exact border rank in AC-03;
no asymptotic result is used as a finite-size certificate.

\item \emph{AC03 round-two verified research package}, supplied prior artifact,
\texttt{AC03\_round2\_research\_package.zip}, report Theorems 4.1, 5.1, 6.2.
Retained unchanged under \texttt{prior/}; its included manifest and combined
verifier are rerun. This is the computational dependency for
Theorems~\ref{thm:prior} and \ref{thm:pair}, not an external peer-reviewed source.
\end{enumerate}
\endgroup
\end{document}
